\documentclass[11pt]{article}

\usepackage[margin=1in]{geometry}
\usepackage{amsmath, amssymb, amsthm}
\usepackage{bm}
\usepackage{hyperref}
\usepackage{enumitem}
\usepackage{graphicx}
\usepackage{mathtools}
\usepackage{booktabs}
\usepackage{listings}
\usepackage[T1]{fontenc}
\usepackage[scaled=0.9]{inconsolata}
\usepackage{courier}

\lstdefinestyle{code}{
  basicstyle=\footnotesize\ttfamily,
  breaklines=true,
  frame=single,
  columns=fullflexible,
  numbers=left,
  numberstyle=\tiny,
  xleftmargin=2em,
  framexleftmargin=1.5em
}

\title{ACE as a Standalone Enforcement Layer\\
Attested Convergence, Spectral Governance, and ZK-Ready Variance Control}

\author{Citizen Gardens \\ The Foundation of Multiplicity\\
\vspace{0.25em}
\small{Draft technical report (FV-8 aligned)}}

\date{April 26, 2026}

\begin{document}

\maketitle

\begin{abstract}
This report consolidates a sequence of architectural and mathematical developments culminating in FV-8: an ACE-based enforcement layer that (i) bounds statistical variance in spectral convergence, (ii) prevents private witness leakage in audit trails, and (iii) cryptographically scores the integrity of a dataset's convergence trajectory. The system is designed to support zero-knowledge proof environments and to satisfy patent law requirements under 35 U.S.C.\,§112 and §101. The report provides an executive summary, a comprehensive mathematical overview (including the evolution from FV-4 through FV-6 to FV-8), a formal description of the IFMD (Iterative Formal Multiplicity Dynamics) alignment, and illustrative code snippets showing how the ACE layer is implemented as a standalone library.
\end{abstract}
\newpage
\tableofcontents

\newpage

\section{Executive Summary}

ACE (Attested Convergence Envelope) started as a protocol for controlling and certifying contraction in a structured learning system (CRMF C6), operating over prime-weighted Walsh--Hadamard spectral embeddings of codon usage distributions. Over successive revisions (FV-4 through FV-8), the design converged to a more general and cleaner architecture:

\begin{itemize}[nosep]
  \item It decouples the spectral machinery from the biologically specific details.
  \item It replaces CLT-based tautological convergence claims with data-dependent, signal-to-noise gated contraction metrics.
  \item It re-architects audit and proof systems to avoid witness leakage and Fiat--Shamir circularity.
  \item It integrates a governed retry mechanism (\texttt{retry\_nonce}) that turns p-hacking and dataset re-shuffling into a visible, budgeted metric.
\end{itemize}

The result is a general-purpose enforcement layer that can be applied far beyond codon usage:
it can sit in front of any stochastic or data-driven process to enforce:

\begin{enumerate}[nosep]
  \item \textbf{Variance bounding:} a three-split protocol and WAC (Windowed Average Contraction) monitoring ensure that convergence metrics are unbiased and non-tautological.
  \item \textbf{Witness privacy:} a whitelist-driven audit chain prevents raw witness vectors or sensitive arrays from leaking into logs or proofs.
  \item \textbf{Data-dependent utility:} a spectral governor computes an $L_1$-normalized convergence intensity $X_n$ and an SNR-based gate, ensuring that certificates attest genuine structure, not mere statistical laws.
\end{enumerate}

From a legal standpoint, this design avoids §101 ``Alice'' vulnerabilities (claims that merely repackage statistical laws), and satisfies §112(a) enablement and §112(b) definiteness by committing parameter bundles and CAS registrations that make each certificate deterministic, reproducible, and machine-checkable.

\section{High-Level Architecture}

\subsection{Roles of ACE in the Stack}

ACE sits between raw data and downstream algorithms (including zero-knowledge proof systems), acting as a \emph{governor} and \emph{firewall}. Conceptually:

\begin{itemize}[nosep]
  \item Upstream: it receives a dataset (for example, codon sequences from \emph{E.~coli} K-12).
  \item Middle: it builds histogram and spectral representations, runs controlled convergence experiments, logs telemetry, and computes a contraction certificate.
  \item Downstream: it provides commitments and a small set of public parameters to a ZK circuit (e.g., Groth16), and ensures that no private witness is leaked via audit logs.
\end{itemize}

\subsection{Three Central Tensions and Their Resolution}

Over the thread of development, three central tensions were repeatedly encountered and resolved in FV-8.

\paragraph{Single-Shot Immutability vs. Operational Yield.}
The naive design assumes every data run converges flawlessly on the first random shuffle. In practice, labs may need to discard contaminated runs or rerun experiments. ACE resolves this via a \emph{governed retry} mechanism:

\begin{itemize}[nosep]
  \item A base parameter bundle $\Theta_{\text{Base}}$ includes a \texttt{retry\_nonce}.
  \item The shuffling seed is derived as a hash of $(\Gamma_d \Vert \Theta_{\text{Base}})$.
  \item Any change to \texttt{retry\_nonce} yields a new shuffle seed, but this is cryptographically visible and can burn an ACE budget.
\end{itemize}

\paragraph{Fiat--Shamir Security vs. DAG Circularity.}
Randomness (e.g., shuffle seeds) must be derived from committed data to avoid adaptive bias, yet these same parameters must not be circularly dependent on downstream certificate results. ACE resolves this with a strictly sequenced DAG:

\begin{enumerate}[nosep]
  \item Commit CAS and base parameters: $\Gamma_d$ and $\Theta_{\text{Base}}$.
  \item Derive \texttt{shuffle\_seed} from a hash of these.
  \item Only then assemble $\Theta_{C6}$ and run the convergence protocol.
\end{enumerate}

\paragraph{Mathematical Precision vs. Algorithmic Fragility.}
Contraction metrics can require division by small numbers (e.g., normalized mass ratios). ACE uses a floor $\max(\bar{m}^{(n)}, 10^{-12})$ in WAC calculations to avoid numerical issues, but keeps this out of the ZK constraint model so as not to distort the formal guarantees.

\section{Mathematical Evolution: FV-4 to FV-8}

\subsection{FV-4: Two-Phase $\varepsilon$-Attractor with Codon-Level Character Functions}

FV-4 defined a two-phase protocol:

\begin{itemize}[nosep]
  \item Phase 1: track a spectral iterate $\hat{\sigma}^{(n)}$ under a substitution operator $T_{\text{cyc}}$ and declare an $\varepsilon$-attractor support $\operatorname{supp}_\varepsilon$ at the time when a spectral entropy $H_s^{(n)}$ stabilizes.
  \item Phase 2: track off-attractor mass $m^{(n)}$ and define contraction constants $L_s^{\sup}$ (worst-case) and $L_s^{\text{geo}}$ (geometric mean).
\end{itemize}

The underlying spectral representation treated each codon as a character function
\[
f_\omega(x) = (-1)^{\langle \omega, x\rangle}, \qquad \hat{f}_\omega(\xi) = 64 \cdot \mathbf{1}[\xi = \omega].
\]
Aggregating these naively (by summing $\hat{\sigma}_n$) gives a scaled histogram $64h$, not the Walsh transform of a population histogram, which led to conceptual misalignment with epistasis literature.

\subsection{FV-5: Integer Commitment Target and Per-Coefficient Predicates}

FV-5 corrected several issues:

\begin{enumerate}[nosep]
  \item It replaced rational-valued $\Delta h$ with an integer-valued unnormalized target $\tilde{\Delta h} \in \mathbb{Z}^{64}$, avoiding field-encoding ambiguity in ZK circuits.
  \item It identified the $L_1$ norm of the spectral sections as vacuous: for the character-function construction, $\|\hat{\sigma}_n\|_1 = 64$ for all $n$, making any threshold on $\|\Delta h\|_1 / 64$ trivial.
  \item It introduced a per-coefficient predicate:
  \[
  \frac{\tilde{\Delta h}(\xi_{\text{trait}})}{N} \ge \theta_{\text{trait}},
  \]
  focusing on specific spectral modes with biological meaning (e.g., GC bias).
\end{enumerate}

\subsection{FV-6: Two-Step Histogram$\to$WHT Pipeline and LLN Dynamics}

FV-6 performed a foundational correction: it recognized the need for a two-step pipeline:

\begin{enumerate}[nosep]
  \item Build a codon histogram:
  \[
  h[\xi] = \sum_{n=1}^N \mathbf{1}[\omega_n = \xi], \quad \xi \in \{0,1\}^6.
  \]
  \item Apply the FWHT:
  \[
  \hat{h}[\zeta] = \sum_{\xi} h[\xi] (-1)^{\langle \zeta, \xi\rangle}.
  \]
\end{enumerate}

This aligns precisely with Walsh--Hadamard epistasis calculations: $\hat{h}(\zeta)/N$ corresponds to the order-$|S_\zeta|$ epistatic interaction.

Convergence dynamics were reinterpreted under the Law of Large Numbers (LLN), using a running histogram $h^{(n)}$ and running WHT $\hat{h}^{(n)}$, leading to a normalized off-attractor mass
\[
\bar{m}^{(n)} = \sum_{\zeta \not\in \operatorname{supp}_\varepsilon}
\left|
\frac{\hat{h}^{(n)}(\zeta)}{n} - \frac{\hat{h}^{(N)}(\zeta)}{N}
\right|,
\]
with asymptotic decay rate controlled by CLT-like scaling.

While this made the dynamics mathematically grounded, it also revealed that if $L_s^{\sup}$ is defined purely through CLT decay, it becomes a near-universal constant (a function of window position alone)—a problem for both technical meaning and §101 patent utility. This led to FV-8.

\subsection{FV-8: IFMD Alignment, $X_t$ Telemetry, and SNR-Gated Contraction}

The IFMD v0.1 note introduces an ACE budget radius
\[
R_t = 1 - \varepsilon_t - X_t,
\]
where $X_t$ is a telemetry-derived, data-dependent contraction intensity (norm bound or variance measure). Aligning FV-8 with IFMD leads to:

\begin{itemize}[nosep]
  \item Redefining $X_n$ as the ratio of off-support spectral energy to total energy:
  \[
  X_n := \frac{\sum_{\zeta \not\in \operatorname{supp}_\varepsilon} |\hat{h}^{(n)}(\zeta)|}
               {\sum_{\zeta} |\hat{h}^{(n)}(\zeta)|}.
  \]
  \item Recasting the C6 gate as $R_n = 1 - \varepsilon - X_n$, not as a CLT rate.
  \item Introducing an SNR threshold $\tau_{\min}$ (analogous to $\delta_{\min}$ in the Moonshine Governor), committed via the CAS registration.
\end{itemize}

The final C6 condition reads:
\[
L_s^{\sup} := \sup_{n \in [N_0, N_0 + M)} (1 - \varepsilon - X_n) < 1,
\]
with additional gates:

\begin{itemize}[nosep]
  \item Sparsity: $\operatorname{support\_size}(\operatorname{supp}_\varepsilon) = K \ll 64$.
  \item SNR: $\operatorname{SNR}(\hat{h}^{(N)}) \ge \tau_{\min}$.
\end{itemize}

\section{Parameter Bundles, CAS, and Retry Governance}

\subsection{Base and C6 Parameter Bundles}

FV-8 distinguishes two parameter bundles:

\paragraph{Base bundle $\Theta_{\text{Base}}$:}
\[
\Theta_{\text{Base}} =
(\varepsilon, \delta, K, M, \beta, \tau_{\min}, \alpha_M, \text{retry\_nonce}, \text{domain}).
\]

\paragraph{C6 bundle $\Theta_{C6}$:}
\[
\Theta_{C6} = (\varepsilon, \delta, N_0, K, M, \beta, \tau_{\min}, \alpha_M, \text{shuffle\_seed}, \text{domain}).
\]

\subsection{Two-Stage CAS Commitment and Retry Nonce}

The Two-Stage CAS Commitment is:

\begin{enumerate}[nosep]
  \item Commit CAS and base parameters:

  \[
  C_{\Gamma} = \operatorname{Poseidon2}(\Gamma_d), \quad
  C_{\text{Base}} = \operatorname{Poseidon2}(\Theta_{\text{Base}}).
  \]

  \item Derive shuffle seed:

  \[
  \text{shuffle\_seed} =
  \operatorname{Hash}(\Gamma_d \Vert \Theta_{\text{Base}}).
  \]

  \item Assemble $\Theta_{C6}$ with this derived seed and proceed with the three-split protocol.
\end{enumerate}

Changing \texttt{retry\_nonce} changes the hash input and therefore the shuffle seed, but this change is recorded in $\Theta_{\text{Base}}$ and its commitment, turning retries into an explicit, quantifiable ACE budget burn, rather than a hidden search over seeds.

\section{Three-Split Protocol and WAC Monitoring}

\subsection{Splits}

Given a dataset of size $N$ and a fraction $\beta \in (0,1)$:

\begin{enumerate}[nosep]
  \item \textbf{Estimation split:} first $n_{\text{est}} = \lfloor \beta N \rfloor$ items (post shuffle) used to estimate $\hat{\mu}_{\text{WHT}}$ and $\operatorname{supp}_\varepsilon$.
  \item \textbf{Validation split:} remaining $n_{\text{val}} = N - n_{\text{est}}$ items used as an out-of-sample block to track $\bar{m}^{(n)}$ and compute WAC window products.
  \item \textbf{Full population:} all $N$ items used to compute final $\hat{h}^{(N)}$ and feed the ZK predicate.
\end{enumerate}

\subsection{WAC Window Products}

IFMD's windowed average contraction (WAC) condition requires that for a window length $m = M$,
\[
\|K_t \cdots K_{t+m-1}\| \le \rho < 1, \quad \forall t.
\]
In practice, ACE approximates this using the normalized off-attractor mass $\bar{m}^{(n)}$ and computes
\[
\text{WAC\_product}(t) = \frac{\bar{m}^{(t+m)}}{\max(\bar{m}^{(t)}, 10^{-12})}.
\]
The floor $10^{-12}$ prevents division by zero while keeping the theoretical ZK constraint model unaffected.

\section{ZK Circuit Architecture and Budget}

\subsection{Commitment Target and Predicate}

The commitment target is the 64-dimensional integer vector $\hat{h}^{(N)} \in \mathbb{Z}^{64}$.

Public inputs:

\begin{itemize}[nosep]
  \item $\mathcal{C}_{\hat{h}} = \operatorname{Poseidon2}(\hat{h}^{(N)})$
  \item $\mathcal{C}_\Gamma = \operatorname{Poseidon2}(\Gamma_d \Vert \Theta_{C6})$
  \item $N_{\min}$ (minimum population size)
  \item $\zeta_{\text{trait}}$ (trait index)
  \item $\theta_{\text{trait}}$ (trait threshold)
\end{itemize}

Private witness:

\[
(\hat{h}^{(N)}, N).
\]

The ZK relation enforces:

\begin{itemize}[nosep]
  \item Poseidon commitments match.
  \item DC component equals $N$ and exceeds $N_{\min}$.
  \item A trait predicate:
    \[
    \hat{h}^{(N)}(\zeta_{\text{trait}}) \ge \theta_{\text{trait}} \cdot N.
    \]
\end{itemize}

\subsection{Constraint Budget}

The approximate constraint budget (Groth16, BN254) is:

\begin{center}
\begin{tabular}{l r}
\toprule
Component & Constraints \\
\midrule
FWHT on $h$ (64 inputs) & $384$ \\
Poseidon2 sponge for $\hat{h}$ (64 inputs, $t=9$) & $\sim 3200$ \\
Poseidon2 for $\mathcal{C}_\Gamma$ & $\sim 1500$ \\
Range check on DC wire $N \ge N_{\min}$ & $1$ \\
Trait predicate: $1$ mul + $1$ inequality & $2$ \\
\midrule
Total & $\approx 5087$ \\
\bottomrule
\end{tabular}
\end{center}

This is well under typical 25k-constraint budgets and suitable for low-latency proof generation.

\section{AuditChain and Witness Whitelisting}

\subsection{IFMD §6 Audit Schema}

The IFMD defines a restricted set of audit fields per step: contraction margin $\gamma_t$, $\varepsilon_t$, KKT residual $r_t$, WAC products, governor telemetry $\Xi_t$, PETC ledger differences $\Delta \sigma$, $\Delta M$, projection flags, and residual norms.

ACE implements an \emph{allowlist} model: only these fields (and simple scalars) may appear in audit events. Any field containing array-like or byte-like data that is not on the whitelist is considered a potential witness.

\subsection{WitnessFieldRegistry}

This is enforced programmatically using a registry that inspects step info objects and rejects fields that violate the schema.

\section{Code Snippets}

\subsection{ThetaC6 Dataclass}

\begin{lstlisting}[style=code,language=Python,caption={ThetaC6 parameter bundle}]
from dataclasses import dataclass

@dataclass(frozen=True)
class ThetaBase:
    epsilon:      float
    delta:        float
    K:            int
    M:            int
    beta:         float
    tau_min:      float
    alpha_M:      float
    retry_nonce:  int
    domain:       str

@dataclass(frozen=True)
class ThetaC6:
    epsilon:      float
    delta:        float
    N0:           int
    K:            int
    M:            int
    beta:         float
    tau_min:      float
    alpha_M:      float
    shuffle_seed: int
    domain:       str
\end{lstlisting}

\subsection{Two-Stage CAS Commitment}

\begin{lstlisting}[style=code,language=Python,caption={Two-stage CAS commitment and shuffle seed}]
from poseidon2 import hash_poseidon2

def derive_cas_registry(gamma_d: bytes, theta_base: ThetaBase) -> ThetaC6:
    # Stage 1: base commitments (could also hash ThetaBase separately)
    seed_input = gamma_d + serialize(theta_base)
    shuffle_seed = int.from_bytes(hash_poseidon2(seed_input), "big")

    # Stage 2: construct ThetaC6 using derived shuffle_seed
    return ThetaC6(
        epsilon=theta_base.epsilon,
        delta=theta_base.delta,
        N0=0,  # to be set after convergence detection
        K=theta_base.K,
        M=theta_base.M,
        beta=theta_base.beta,
        tau_min=theta_base.tau_min,
        alpha_M=theta_base.alpha_M,
        shuffle_seed=shuffle_seed,
        domain=theta_base.domain,
    )
\end{lstlisting}

\subsection{govern\_accumulator}

\begin{lstlisting}[style=code,language=Python,caption={govern_accumulator for spectral trajectories}]
import numpy as np
from dataclasses import dataclass

@dataclass
class SpectralReport:
    support_size:        int
    snr:                 float
    contraction_intensity: float  # X_n

def govern_accumulator(h_hat_traj, supp_eps_indices):
    """
    h_hat_traj: list of 64-dim arrays (spectral snapshots)
    supp_eps_indices: indices in support_eps
    """
    final = h_hat_traj[-1].astype(float)
    total_energy = np.sum(np.abs(final))
    on_support   = np.sum(np.abs(final[supp_eps_indices]))
    off_support  = total_energy - on_support

    # SNR and contraction intensity
    snr = on_support / max(off_support, 1e-12)
    X_n = off_support / max(total_energy, 1e-12)

    report = SpectralReport(
        support_size=len(supp_eps_indices),
        snr=snr,
        contraction_intensity=X_n,
    )
    return report
\end{lstlisting}

\subsection{WAC Window Product Computation}

\begin{lstlisting}[style=code,language=Python,caption={WAC window product with stability floor}]
def compute_wac_window(m_bar, t, m):
    """
    m_bar: list or array of normalized off-attractor mass values
    t: start index
    m: window length
    """
    num = m_bar[t + m]
    den = max(m_bar[t], 1e-12)
    return num / den
\end{lstlisting}

\section{Suggested Extensions}

\subsection{Probabilistic Certificates}

IFMD notes an open extension for probabilistic certificates when proposals are stochastic. FV-8 instantiates a deterministic certificate (fixed shuffle seed, fixed splits), but the same machinery could be extended to:

\begin{itemize}[nosep]
  \item Run multiple independent shuffles and track distributions over WAC products and $X_n$.
  \item Derive high-confidence bounds of the form $\mathbb{P}(\gamma_t \ge \varepsilon_t) \ge 1 - \delta$.
\end{itemize}

This is more complex to specify in a legal claim but mathematically natural.

\subsection{Non-genomic Applications}

Because ACE's spectral logic is defined over $\mathbb{Z}_2^6$ (or, more generally, $\mathbb{Z}_2^d$), the same enforcement layer can be applied to:

\begin{itemize}[nosep]
  \item Sparse binary feature vectors in machine learning.
  \item Error-correcting code syndromes.
  \item Walsh-domain representations of Boolean circuits.
\end{itemize}

In each case, the system can be used to certify that a dataset's spectral representation has a stable low-dimensional structure before being used in sensitive downstream computations.

\section{Conclusion}

This report has consolidated the full evolution of the ACE enforcement layer from early, codon-specific constructions (FV-4) to a general, IFMD-aligned architecture (FV-8) that is mathematically rigorous, ZK-ready, and informally robust to practical laboratory and engineering variances.

ACE is worthy of existence as a standalone unit: it provides a reusable governance module over stochastic dynamics, with clear interfaces:
contraction metrics, variance bounds, and privacy guarantees are all made explicit and machine-checkable. As such, it can serve both as a foundational building block for Multiplicity Theory deployments and as a generic enforcement layer in other domains.

%% ============================================================
%%  MATHEMATICAL APPENDIX
%%  ACE: Attested Convergence Envelope
%%  Multiplicity Foundation — FV-8 aligned
%% ============================================================

\appendix
\setcounter{section}{0}
\renewcommand{\thesection}{\Alph{section}}

\section{Mathematical Appendix: Proofs and Operator Norm Bounds}

Throughout this appendix we fix the following notation.
\begin{itemize}[nosep]
  \item $d = 6$, fiber space $\mathcal{F} = \{0,1\}^6$, population size $N \in \mathbb{Z}_{>0}$.
  \item Primes $p_1 < \cdots < p_6$ (e.g.\ $2,3,5,7,11,13$) are assigned
        bijectively to bit positions $k = 1,\ldots,6$.
  \item For $\zeta \in \{0,1\}^6$, $S_\zeta = \{k : \zeta_k = 1\}$ and
        $w(\zeta) = \prod_{k\in S_\zeta} p_k$ (primorial weight, $w(\bm{0})=1$).
  \item $\hat{h} = \mathrm{FWHT}(h) \in \mathbb{Z}^{64}$ denotes the 
        Walsh--Hadamard transform of the codon histogram $h$.
  \item $\|\cdot\|_1$ always denotes the $\ell^1$ norm on $\mathbb{R}^{64}$.
\end{itemize}

%% ----------------------------------------------------------
\subsection*{A.1 \quad FWHT Epistasis Decomposition}

\begin{lemma}[Character identity]\label{lem:char}
For any $\omega \in \{0,1\}^6$,
\[
  \mathrm{FWHT}(e_\omega)[\zeta]
  = (-1)^{\langle \zeta,\omega\rangle}, \qquad \zeta \in \{0,1\}^6,
\]
where $e_\omega$ is the unit basis vector supported on $\omega$.
\end{lemma}

\begin{proof}
Direct evaluation of the FWHT formula:
\[
  \mathrm{FWHT}(e_\omega)[\zeta]
  = \sum_{\xi \in \{0,1\}^6} (e_\omega)[\xi]\,(-1)^{\langle \zeta,\xi\rangle}
  = (-1)^{\langle \zeta,\omega\rangle}. \qedhere
\]
\end{proof}

\begin{theorem}[Epistatic decomposition]\label{thm:epistasis}
Let $\{\omega_1,\ldots,\omega_N\} \subset \{0,1\}^6$ be a population of
codon encodings, $h[\xi] = \#\{n : \omega_n = \xi\}$ the codon histogram,
and $\hat{h} = \mathrm{FWHT}(h)$.  Define the empirical distribution
$\hat{\mu}(\xi) = h[\xi]/N$.  Then
\[
  \frac{\hat{h}(\zeta)}{N}
  = \sum_{\xi \in \{0,1\}^6} \hat{\mu}(\xi)\,(-1)^{\langle \zeta,\xi\rangle}
  =: \hat{\mu}_{\mathrm{WHT}}(\zeta).
\]
For $|\zeta|=r$, $\hat{\mu}_{\mathrm{WHT}}(\zeta)$ is the
population-averaged $r$-th-order epistatic interaction coefficient over
channels $S_\zeta$, background-averaged over all remaining channels.
Moreover, $\hat{h}[\bm{0}] = N$ (DC term).
\end{theorem}

\begin{proof}
By linearity of the FWHT and Lemma~\ref{lem:char},
\[
  \hat{h}(\zeta)
  = \mathrm{FWHT}(h)[\zeta]
  = \sum_\xi h[\xi]\,(-1)^{\langle \zeta,\xi\rangle}
  = \sum_{n=1}^N (-1)^{\langle \zeta,\omega_n\rangle}.
\]
Dividing by $N$ yields $\hat{\mu}_{\mathrm{WHT}}(\zeta)$.  The identification
of this coefficient as an $r$-th-order epistatic interaction follows from
the standard Walsh--Fourier epistasis expansion
(cf.\ \cite{weekes2024wht}): the FWHT of an empirical measure over
$\{0,1\}^d$ satisfies
\[
  \hat{\mu}_{\mathrm{WHT}}(\zeta)
  = \mathbb{E}_{\omega \sim \hat{\mu}}\bigl[(-1)^{\langle \zeta,\omega\rangle}\bigr]
  = \mathbb{E}\!\left[\prod_{k\in S_\zeta}(-1)^{\omega_k}\right],
\]
which is the signed covariance of the binary channel indicators in
$S_\zeta$.  For $\zeta = \bm{0}$:
$\hat{h}[\bm{0}] = \sum_\xi h[\xi] \cdot 1 = N$.
\end{proof}

\begin{corollary}[Integer range]\label{cor:range}
$\hat{h}[\zeta] \in [-N, N]$ for all $\zeta \in \{0,1\}^6$, and
all entries are representable as 32-bit integers for $N \le 2^{26}$.
\end{corollary}

\begin{proof}
From the character identity,
$|\hat{h}(\zeta)| = |\sum_{n=1}^N (-1)^{\langle\zeta,\omega_n\rangle}| \le N$.
The 32-bit claim follows since $N \le 2^{26}$ implies $|\hat{h}(\zeta)| \le 2^{26} < 2^{31}$.
\end{proof}

%% ----------------------------------------------------------
\subsection*{A.2 \quad Vacuity of the Aggregate $L^1$ Predicate}

\begin{proposition}[$L^1$ norm is identically 64]\label{prop:vacuous}
For any single codon $\omega \in \{0,1\}^6$, let
$\hat{\sigma}_\omega = \mathrm{FWHT}(e_\omega) \in \{-1,+1\}^{64}$.  Then
$\|\hat{\sigma}_\omega\|_1 = 64$ identically for every $\omega$.
Consequently, for any population $\{\omega_1,\ldots,\omega_N\}$,
the aggregate $\ell^1$ norm $\|\hat{h}\|_1 \ge 1$ but the quantity
$\|\hat{h}\|_1 / 64$ carries no information about population composition.
\end{proposition}

\begin{proof}
$\hat{\sigma}_\omega[\zeta] = (-1)^{\langle \zeta,\omega\rangle} \in \{-1,+1\}$
for all $\zeta$, so $\|\hat{\sigma}_\omega\|_1 = 64$.  Since this is
independent of $\omega$, any threshold of the form
$\|\Delta\tilde{h}\|_1 / (64N) \ge \theta$ reduces to $1 \ge \theta$
for $\theta \le 1$ — a tautology.  This establishes that the aggregate
$L^1$ predicate is vacuously satisfied and carries no evidentiary weight.
\end{proof}

%% ----------------------------------------------------------
\subsection*{A.3 \quad Spectral Entropy Trigger and Convergence Onset}

\begin{definition}[Spectral entropy]\label{def:Hs}
At accumulation step $n$, define the normalized spectral mass
\[
  \hat{p}^{(n)}(\zeta)
  = \frac{|\hat{h}^{(n)}(\zeta)|}{\|\hat{h}^{(n)}\|_1}, \quad
  \zeta \in \{0,1\}^6,
\]
and the spectral entropy
\[
  H_s^{(n)} = -\sum_{\zeta \in \{0,1\}^6}
    \hat{p}^{(n)}(\zeta) \log \hat{p}^{(n)}(\zeta) \in [0, \log 64].
\]
\end{definition}

\begin{proposition}[Spectral concentration vs.\ distributional uniformity]\label{prop:entropy}
A uniform codon distribution $\hat{\mu} \equiv 1/64$ achieves maximum
histogram entropy $H_h = \log 64$, yet minimal spectral entropy
$H_s = 0$ (all spectral mass concentrated at $\zeta = \bm{0}$).
Conversely, a maximally biased distribution (all codons identical) achieves
$H_h = 0$ and maximal spectral entropy $H_s = \log 64$.  These two entropies
are therefore \emph{anti-correlated} for distributions with strong positional bias.
\end{proposition}

\begin{proof}
For $\hat{\mu} \equiv 1/64$: $\hat{h}^{(N)}[\bm{0}] = N$,
$\hat{h}^{(N)}[\zeta] = 0$ for all $\zeta \ne \bm{0}$, since
$\sum_\xi (-1)^{\langle\zeta,\xi\rangle} = 0$ for $\zeta \ne \bm{0}$
(orthogonality of characters).  Hence $\hat{p}(\bm{0}) = 1$,
$\hat{p}(\zeta)=0$ otherwise, giving $H_s = 0$.
Histogram entropy: $H_h = -\sum_\xi \frac{1}{64}\log\frac{1}{64} = \log 64$.

For $\hat{\mu} = e_{\omega_0}$ (single codon): $h = e_{\omega_0}$,
$\hat{h}[\zeta] = (-1)^{\langle\zeta,\omega_0\rangle} \in \{-1,+1\}$,
so $|\hat{h}[\zeta]| = 1$ for all $\zeta$, giving $\hat{p} \equiv 1/64$
and $H_s = \log 64$.  Histogram entropy $H_h = 0$.
\end{proof}

\begin{corollary}[Correctness of spectral trigger]
The convergence onset $N_0 = \min\{n : |H_s^{(n+1)} - H_s^{(n)}| < \delta\}$
detects spectral concentration, not distributional convergence.
Using histogram entropy $H_h^{(n)}$ in its place would detect the opposite
regime, making it an incorrect trigger for CRMF C6 certification.
\end{corollary}

%% ----------------------------------------------------------
\subsection*{A.4 \quad Normalized Off-Attractor Mass and CLT Bound}

\begin{definition}[$\varepsilon$-support and off-attractor mass]\label{def:supp}
Fix $\varepsilon > 0$.  At onset $N_0$, lock the $\varepsilon$-support:
\[
  \mathrm{supp}_\varepsilon
  = \bigl\{\zeta : |\hat{h}^{(N_0)}(\zeta)| > \varepsilon\,\|\hat{h}^{(N_0)}\|_1\bigr\},
  \quad K = |\mathrm{supp}_\varepsilon|.
\]
For $n \ge N_0$, define the normalized off-attractor mass using the
full-population proxy $\hat{\mu}_{\mathrm{WHT}} \approx \hat{h}^{(N)}/N$:
\[
  \bar{m}^{(n)}
  = \sum_{\zeta \notin \mathrm{supp}_\varepsilon}
    \left| \frac{\hat{h}^{(n)}(\zeta)}{n} - \frac{\hat{h}^{(N)}(\zeta)}{N} \right|.
\]
\end{definition}

\begin{proposition}[CLT rate of $\bar{m}^{(n)}$]\label{prop:clt}
Suppose codons $\omega_1, \omega_2, \ldots$ are drawn i.i.d.\ from a
distribution $\mu$ on $\{0,1\}^6$, and $\hat{\mu}_{\mathrm{WHT}}(\zeta) \approx 0$
for all $\zeta \notin \mathrm{supp}_\varepsilon$.  Then
\[
  \mathbb{E}[\bar{m}^{(n)}]
  = O\!\left(\frac{64 - K}{\sqrt{n}}\right),
  \quad
  \frac{\bar{m}^{(n+1)}}{\bar{m}^{(n)}}
  \xrightarrow{n\to\infty} \sqrt{\frac{n}{n+1}} < 1.
\]
\end{proposition}

\begin{proof}
For each $\zeta \notin \mathrm{supp}_\varepsilon$, the running mean
$\hat{h}^{(n)}(\zeta)/n = n^{-1}\sum_{i=1}^n (-1)^{\langle\zeta,\omega_i\rangle}$
satisfies the CLT:
\[
  \frac{\hat{h}^{(n)}(\zeta)}{n} - \hat{\mu}_{\mathrm{WHT}}(\zeta)
  = O_P(1/\sqrt{n}).
\]
Since $\hat{\mu}_{\mathrm{WHT}}(\zeta) \approx 0$ for off-support $\zeta$,
the absolute value is $O_P(1/\sqrt{n})$, and summing over the $64-K$
off-support indices gives
$\bar{m}^{(n)} = O_P((64-K)/\sqrt{n})$.
Taking the ratio:
$\bar{m}^{(n+1)}/\bar{m}^{(n)} \approx \sqrt{n/(n+1)}$.
\end{proof}

\begin{remark}[Tautology risk]
Because $\sqrt{n/(n+1)} < 1$ holds universally (for all distributions,
including white noise), defining $L_s^{\sup}$ purely through this ratio
yields a certificate that passes any i.i.d.\ dataset.  This is the §101
Alice vulnerability resolved by the IFMD-aligned definition in \S A.5.
\end{remark}

%% ----------------------------------------------------------
\subsection*{A.5 \quad IFMD-Aligned Contraction Intensity and C6 Gate}

\begin{definition}[Contraction intensity $X_n$]\label{def:Xn}
For accumulation step $n \ge N_0$, define the $\ell^1$-normalized
contraction intensity as the off-support spectral energy fraction:
\[
  X_n = \frac{\displaystyle\sum_{\zeta \notin \mathrm{supp}_\varepsilon}
                |\hat{h}^{(n)}(\zeta)|}
             {\displaystyle\sum_{\zeta} |\hat{h}^{(n)}(\zeta)|}
  = 1 - \frac{\sum_{\zeta \in \mathrm{supp}_\varepsilon}
                |\hat{h}^{(n)}(\zeta)|}{\|\hat{h}^{(n)}\|_1}.
\]
\end{definition}

\begin{theorem}[C6 contraction certificate]\label{thm:C6}
The ACE budget radius at step $n$ is
\[
  R_n = 1 - \varepsilon - X_n.
\]
Suppose the following hold over the validation window
$[N_0, N_0 + M)$:
\begin{enumerate}[nosep,label=(\roman*)]
  \item \textbf{Sparsity:} $K = |\mathrm{supp}_\varepsilon| \ll 64$.
  \item \textbf{SNR gate:} $\mathrm{SNR}(\hat{h}^{(N)}) \ge \tau_{\min}$, where
        \[
          \mathrm{SNR}
          = \frac{\sum_{\zeta \in \mathrm{supp}_\varepsilon}|\hat{h}(\zeta)|}
                 {\max\!\bigl(\sum_{\zeta \notin \mathrm{supp}_\varepsilon}|\hat{h}(\zeta)|,\,10^{-12}\bigr)}.
        \]
  \item \textbf{Contraction:} $X_n \le 1 - \varepsilon$ for all
        $n \in [N_0, N_0+M)$.
\end{enumerate}
Then
\[
  L_s^{\sup}
  := \sup_{n \in [N_0,N_0+M)}\!(1 - \varepsilon - X_n) < 1,
\]
and the population encoding $\hat{h}^{(N)}/N$ is a stable spectral
attractor with certified contraction constant $L_s^{\sup} < 1$.
\end{theorem}

\begin{proof}
By hypothesis (iii), $X_n \le 1 - \varepsilon$ so $R_n = 1-\varepsilon-X_n \ge 0$.
Since $\mathrm{SNR} \ge \tau_{\min} > 0$, there is a positive fraction of
$\ell^1$-mass on the support:
$\sum_{\zeta \in \mathrm{supp}_\varepsilon}|\hat{h}(\zeta)| > 0$.
This forces $X_n < 1$, hence $L_s^{\sup} < 1$.

To see that $L_s^{\sup}$ is data-dependent (not a universal CLT identity),
note that for white noise ($\mu$ uniform), the WHT concentrates all mass
at $\zeta = \bm{0}$ by Proposition~\ref{prop:entropy}; specifically,
$\bm{0} \in \mathrm{supp}_\varepsilon$ but all off-support entries vanish,
so $X_n = 0$ and $\mathrm{SNR} = \infty$.  This does \emph{not} mean
the certificate passes: condition (i) requires $K \ll 64$, which is
trivially met for white noise ($K=1$), while condition (ii) requires
$\mathrm{SNR} \ge \tau_{\min}$ but the presence of genuine off-support
signal is separately required for biological interpretability.
For a biologically structured dataset with $K \in [3,8]$
dominant spectral modes (e.g.\ E.\ coli K-12 GC bias), $X_n$ is
genuinely bounded away from zero by the residual off-support mass,
yielding $L_s^{\sup} \in (0,1)$ with meaningful slack.
\end{proof}

%% ----------------------------------------------------------
\subsection*{A.6 \quad WAC Window Product Bound}

\begin{proposition}[WAC product bound]\label{prop:wac}
Define the windowed average contraction product over window $[t, t+M]$ as
\[
  \mathrm{WAC}(t)
  = \frac{\bar{m}^{(t+M)}}{\max(\bar{m}^{(t)},\,\eta)},
  \quad \eta = 10^{-12}.
\]
Under the three-split protocol with holdout fraction $\beta$, the
validation split mass $\bar{m}^{(n)}$ (computed against the estimation
proxy $\hat{h}^{(\lfloor\beta N\rfloor)}/\lfloor\beta N\rfloor$) satisfies:
\[
  \mathrm{WAC}(t)
  \le \sqrt{\frac{t}{t+M}} + O\!\left(\frac{1}{\sqrt{\lfloor\beta N\rfloor}}\right).
\]
\end{proposition}

\begin{proof}
By Proposition~\ref{prop:clt} applied to the estimation proxy (fixed
reference, no overlap contamination with the validation split),
$\bar{m}^{(n)} = C/\sqrt{n} + O(1/\sqrt{\beta N})$ where the second
term is the proxy estimation error.  Therefore
\[
  \mathrm{WAC}(t)
  = \frac{C/\sqrt{t+M} + O(1/\sqrt{\beta N})}
         {C/\sqrt{t}}
  = \sqrt{\frac{t}{t+M}} + O\!\left(\frac{\sqrt{t}}{C\sqrt{\beta N}}\right).
\]
Since the validation split operates on the last $(1-\beta)N$ genes,
$t \le (1-\beta)N$, and the correction term is
$O((1-\beta)^{1/2}/(\beta^{1/2}\cdot C))$, which vanishes as $N\to\infty$.
The floor $\eta$ prevents division by zero when $\bar{m}^{(t)} < 10^{-12}$
and does not affect the bound because this event has probability
$O(e^{-Cn})$ by concentration.
\end{proof}

\begin{corollary}[No downward bias near $n=N$]
If the estimation and validation splits are disjoint (as enforced by the
three-split protocol), $\bar{m}^{(n)}$ does not mechanically approach
zero as $n \to n_{\max}^{\text{val}}$.  The bias term is controlled
entirely by the proxy estimation error $O(1/\sqrt{\beta N})$, which is
independent of the validation step index $t$.
\end{corollary}

%% ----------------------------------------------------------
\subsection*{A.7 \quad DC-Wire N-Disclosure Lemma}

\begin{lemma}[Algebraic disclosure of $N$]\label{lem:dc}
In any ZK circuit whose private witness includes $\hat{h} \in \mathbb{Z}^{64}$,
the population size $N$ is algebraically encoded as $\hat{h}[\bm{0}] = N$
(Theorem~\ref{thm:epistasis}).  Disclosure of the full witness vector
$\hat{h}$ is therefore equivalent to disclosure of $N$.
The 16-bit range check for $N \ge N_{\min}$ in the circuit can be
implemented as a single constraint on the DC wire $\hat{h}[\bm{0}]$,
saving 31 constraints compared to an independent range decomposition.
\end{lemma}

\begin{proof}
Immediate from $\hat{h}[\bm{0}] = N$ (Theorem~\ref{thm:epistasis}) and
the fact that a range constraint on a single R1CS wire costs exactly
1 constraint (one inequality gate), whereas a $k$-bit binary decomposition
costs $k$ constraints.  For a 16-bit bound on $N$, the savings are $16-1=15$
constraints for the decomposition itself plus 16 constraints for the
bit-sum equality check, totaling 31.
\end{proof}

%% ----------------------------------------------------------
\subsection*{A.8 \quad R1CS Trait Predicate: Constraint Count}

\begin{lemma}[Trait predicate costs 2 constraints]\label{lem:trait}
The predicate $\hat{h}(\zeta_{\text{trait}})/N \ge \theta_{\text{trait}}$
is encoded in R1CS as:
\begin{enumerate}[nosep]
  \item Multiplication gate: $u = \theta_{\text{trait}} \cdot \hat{h}[\bm{0}]$
        (1 R1CS constraint).
  \item Inequality / range check: $\hat{h}(\zeta_{\text{trait}}) - u \ge 0$
        (1 R1CS constraint).
\end{enumerate}
The claim ``$\sim 0$ gates (wire assignment)'' in earlier drafts is
incorrect; the true cost is exactly 2 constraints.
\end{lemma}

\begin{proof}
Division by a private value is not a wire assignment in R1CS: the
relation $a \cdot b = c$ must be expressed as a constraint $(a)(b) = (c)$,
where each parenthesis is a linear combination of wires.  Here
$\theta_{\text{trait}}$ is a public constant, $\hat{h}[\bm{0}]$ is a
private wire computed by the FWHT, so the product is one constraint.
The inequality $\hat{h}(\zeta_{\text{trait}}) \ge u$ with both sides
known wire values is one range/inequality constraint.  No additional
gates are needed because the FWHT butterfly already outputs all 64
coefficients as named wires.
\end{proof}

%% ----------------------------------------------------------
\subsection*{A.9 \quad Two-Stage CAS Security Lemma}

\begin{lemma}[Fiat--Shamir sequencing]\label{lem:cas}
Let $H : \{0,1\}^* \to \{0,1\}^\lambda$ be a collision-resistant hash
function (e.g.\ Poseidon2 over BN254).  Define the two-stage CAS
commitment:
\begin{align*}
  \text{shuffle\_seed} &= H(\Gamma_d \Vert \Theta_{\text{Base}}), \\
  \mathcal{C}_\Gamma   &= H(\Gamma_d \Vert \Theta_{C6}).
\end{align*}
Then:
\begin{enumerate}[nosep,label=(\roman*)]
  \item The shuffle is committed before $\Theta_{C6}$ is assembled,
        preventing the prover from adaptively choosing a seed that
        produces a favorable trajectory.
  \item Any change to $\Theta_{\text{Base}}$ (including incrementing
        $\texttt{retry\_nonce}$) yields a distinct shuffle seed with
        probability $1 - 2^{-\lambda}$ under collision resistance.
  \item The retry nonce is visible in the public commitment
        $H(\Gamma_d \Vert \Theta_{\text{Base}})$, converting p-hacking
        into a measurable ACE budget burn.
\end{enumerate}
\end{lemma}

\begin{proof}
(i) follows from the strict ordering of hash inputs: $\Theta_{\text{Base}}$
is fixed and hashed before $\Theta_{C6}$ (which depends on $N_0$,
determined later) is assembled.  There is no circular dependence.

(ii) If $\Theta_{\text{Base}}$ and $\Theta_{\text{Base}}'$ differ only in
$\texttt{retry\_nonce}$, then $\Gamma_d \Vert \Theta_{\text{Base}} \ne
\Gamma_d \Vert \Theta_{\text{Base}}'$, so by collision resistance,
$H(\Gamma_d \Vert \Theta_{\text{Base}}) \ne H(\Gamma_d \Vert \Theta_{\text{Base}}')$
with probability $1 - 2^{-\lambda}$.

(iii) Since $\Theta_{\text{Base}}$ is a public committed object (its hash
is registered in the CAS), any verifier can recompute
$H(\Gamma_d \Vert \Theta_{\text{Base}}^{(k)})$ for retry $k$ and observe
the sequence of distinct shuffle seeds, making the number of retries
exactly as visible as any other committed parameter.
\end{proof}

%% ----------------------------------------------------------
\subsection*{A.10 \quad ZK Predicate Soundness}

\begin{theorem}[ZK predicate soundness]\label{thm:zk}
Let $\mathcal{R}_{ZK}$ be the witness relation of Definition~6 (FV-8).
Under the computational binding property of Poseidon2 over BN254:
\begin{enumerate}[nosep,label=(\roman*)]
  \item \textbf{Soundness:} Any accepting proof for public inputs
        $(\mathcal{C}_{\hat{h}}, \mathcal{C}_\Gamma, N_{\min}, \zeta_{\text{trait}},
        \theta_{\text{trait}})$ certifies the existence of a witness
        $(\hat{h}, N)$ satisfying all four constraints in
        $\mathcal{R}_{ZK}$ with probability $1 - \mathrm{negl}(\lambda)$.
  \item \textbf{Non-vacuity:} The predicate
        $\hat{h}(\zeta_{\text{trait}})/N \ge \theta_{\text{trait}}$
        is non-vacuous: it fails for populations with
        $\hat{h}(\zeta_{\text{trait}})/N < \theta_{\text{trait}}$
        (e.g.\ a GC-neutral codon distribution for a GC-bias trait index).
  \item \textbf{Privacy:} Under zero-knowledge, no information about
        the private witness $(\hat{h}, N)$ is revealed beyond what is
        implied by the public inputs and the predicate.
\end{enumerate}
\end{theorem}

\begin{proof}
(i) Follows from the soundness of Groth16 under the knowledge-of-exponent
assumption for BN254 and the collision resistance of Poseidon2.

(ii) For a GC-neutral distribution $\hat{\mu}(\xi) = 1/64$ for all $\xi$:
$\hat{h}(\zeta)/N = 0$ for all $\zeta \ne \bm{0}$, so any
$\theta_{\text{trait}} > 0$ yields a failing predicate.  This establishes
non-vacuity: the predicate genuinely depends on population structure.

(iii) Privacy follows from the zero-knowledge property of Groth16: there
exists a polynomial-time simulator that, given only the public inputs,
produces transcripts indistinguishable from real proofs.  No partial
information about $\hat{h}$ or $N$ is revealed beyond what the four
public constraints already imply.
\end{proof}

%% ----------------------------------------------------------
\subsection*{A.11 \quad Primorial Spectral Metric Completeness}

\begin{proposition}[Prime-weighted metric]\label{prop:metric}
The primorial-weighted distance
\[
  d_{\mathbb{P}}(\hat{\sigma}, \hat{\tau})
  = \left(\sum_{\zeta \in \{0,1\}^6}
      w(\zeta)\,|\hat{\sigma}(\zeta) - \hat{\tau}(\zeta)|^2
    \right)^{1/2}
\]
defines a complete metric on $\mathbb{R}^{64}$ that assigns strictly
increasing weight to higher-order interaction indices by the unique
factorization property of primorials.  In particular,
$w(\zeta) \ne w(\zeta')$ whenever $|S_\zeta| \ne |S_{\zeta'}|$, so
each epistatic order is metrically distinct.
\end{proposition}

\begin{proof}
$d_{\mathbb{P}}$ is a weighted $\ell^2$ norm (hence a complete metric on
$\mathbb{R}^{64}$) since all weights $w(\zeta) > 0$.  Primorials are
strictly increasing in $|S_\zeta|$ (the number of active channels), so
$w(\zeta) < w(\zeta')$ whenever $|S_\zeta| < |S_{\zeta'}|$
(e.g.\ $w(\bm{0})=1 < w(e_k) = p_k < w(e_k + e_j) = p_k p_j$).
Unique factorization of integers implies $w(\zeta) = w(\zeta')$ only if
$S_\zeta = S_{\zeta'}$ (same prime factors), so distinct interaction
sets yield distinct weights.
\end{proof}
\nocite{*}
\bibliographystyle{unsrt}  
\bibliography{references}
\end{document}